%==============================================================================
\section{Deep Learning --- Học sâu và học máy}
%==============================================================================

\begin{ynghia}[title={Nguồn}]
\tsurl{https://www.tutorialspoint.com/machine_learning/deep_machine_learning.htm}.
\end{ynghia}

\begin{dongco}[title={Hai thuật ngữ hay bị dùng lẫn}]
Trong thế giới AI, ``machine learning'' và ``deep learning'' thường được dùng như nhau nhưng không trùng nghĩa.
ML là tập con của AI cho phép máy học mà không lập trình tường minh; deep learning là tập con của ML dùng mạng nơ-ron để xử lý dữ liệu phức tạp.
Bài này so sánh ML và deep learning và mối quan hệ giữa chúng.
\end{dongco}

\subsection{Học máy là gì?}

\begin{dinhnghia}[title={Machine Learning (ML)}]
Học máy (ML) là lĩnh vực con của AI, tự động cho phép máy học từ kinh nghiệm.
Trong ML, phát triển thuật toán là công việc cốt lõi: thuật toán được huấn luyện trên dữ liệu để học pattern ẩn và dự đoán dựa trên những gì đã học; toàn bộ quá trình huấn luyện đôi khi gọi là \textbf{model building}.
\end{dinhnghia}

\begin{ynghia}[title={``Kinh nghiệm'' nghĩa là gì?}]
Khi nói máy học từ kinh nghiệm, ``kinh nghiệm'' ở đây chính là \textbf{huấn luyện thuật toán bằng dữ liệu} --- tương tự con người học từ trải nghiệm, máy học từ dữ liệu khi ta huấn luyện.
ML là một \textbf{kỹ thuật} triển khai giải pháp cho các bài toán liên quan AI; có nhiều cách triển khai, ML là một trong số đó.
\end{ynghia}

\begin{phuongphap}[title={Bốn hướng tiếp cận học từ dữ liệu}]
\begin{itemize}
  \item Supervised learning
  \item Unsupervised learning
  \item Semi-supervised learning
  \item Reinforcement learning
\end{itemize}
Supervised learning huấn luyện trên dữ liệu có nhãn, phù hợp phân loại và hồi quy; ví dụ thuật toán: linear regression, k-nearest neighbors, random forests, \ldots
\end{phuongphap}

\begin{ynghia}[title={Mạng nơ-ron và độ chính xác}]
Nguồn nêu: với mạng nơ-ron, ML đạt độ chính xác cao nhất trong nhiều bài toán.
Mạng nơ-ron có thể xem là một hướng supervised learning phức tạp.
Deep learning là cách khác để triển khai giải pháp ML, dùng mạng nơ-ron học quan hệ phức tạp trong dữ liệu.
\end{ynghia}

\subsection{Học sâu (Deep Learning) là gì?}

\begin{dinhnghia}[title={Deep Learning}]
Deep learning là một kiểu học máy dùng mạng nơ-ron xử lý dữ liệu phức tạp: máy tự học pattern và quan hệ trong dữ liệu bằng nhiều lớp node/nơ-ron liên kết.
Thuật toán deep learning phát hiện và học từ pattern để dự đoán hoặc ra quyết định.
\end{dinhnghia}

\begin{ynghia}[title={Phù hợp với dữ liệu phức tạp}]
Deep learning đặc biệt phù hợp nhận dạng ảnh, giọng nói, NLP, xe tự lái; có thể xử lý lượng dữ liệu rất lớn và học pattern/quan hệ phức tạp.
\end{ynghia}

\begin{vidu}[title={Ví dụ deep learning}]
Nhận diện khuôn mặt, nhận dạng giọng nói, xe tự lái.
\end{vidu}

\subsection{Khác biệt giữa Machine Learning và Deep Learning}

\begin{vidu}[title={Bảng so sánh (dịch từ nguồn)}]
\begin{tabularx}{\linewidth}{>{\raggedright\arraybackslash}p{2.6cm}XX}
\textbf{Tiêu chí} & \textbf{Machine Learning} & \textbf{Deep Learning} \\
\midrule
Định nghĩa &
Tập con AI: máy học không lập trình tường minh; thuật toán học pattern ẩn từ dữ liệu để dự đoán/quyết định trên dữ liệu mới. &
Tập con ML: dùng mạng nơ-ron xử lý dữ liệu phức tạp. \\
Độ phức tạp &
Phương pháp đơn giản hơn: decision tree, linear regression, \ldots &
Phương pháp phức tạp trong mạng nơ-ron (CNN, RNN, GAN, \ldots). \\
Lượng dữ liệu &
Cần nhiều dữ liệu; vẫn hữu ích khi dữ liệu ít. &
Cần nhiều dữ liệu hơn ML; độ chính xác tăng khi dữ liệu tăng. \\
Phương pháp huấn luyện &
Bốn hướng: supervised, unsupervised, semi-supervised, reinforcement. &
Phương pháp phức tạp: CNN, RNN, GAN, \ldots \\
Phụ thuộc phần cứng &
Phương pháp đơn giản hơn $\Rightarrow$ ít lưu trữ và tính toán hơn. &
Phức tạp hơn, dữ liệu lớn hơn $\Rightarrow$ cần nhiều lưu trữ và sức mạnh tính toán hơn. \\
Feature engineering &
Thường phải làm feature engineering \textbf{thủ công}. &
Mô hình deep learning có thể tự thực hiện phần lớn feature engineering. \\
Cách giải bài toán &
Tiếp cận chuẩn; dùng thống kê và toán. &
Dùng thống kê, toán \textbf{kèm} kiến trúc mạng nơ-ron. \\
Thời gian chạy &
Thuật toán ML thường nhanh hơn. &
Huấn luyện lâu hơn vì nhiều tham số và dữ liệu phức tạp. \\
Phù hợp nhất &
Dữ liệu có cấu trúc (structured). &
Dữ liệu phức tạp, không cấu trúc (unstructured). \\
\end{tabularx}
\end{vidu}

\subsection{So sánh sâu hơn}

\begin{ynghia}[title={Bốn điểm (theo nguồn)}]
\begin{enumerate}
  \item ML là nhóm rộng gồm deep learning; deep learning là \textbf{một loại} thuật toán ML dùng mạng nơ-ron.
  \item Thuật toán ML cải thiện theo dữ liệu; deep learning được thiết kế để xử lý dữ liệu phức tạp và nhận pattern mà ML cổ điển có thể bỏ sót.
  \item Deep learning cần nhiều dữ liệu và phần cứng mạnh (GPU) hơn nhiều thuật toán ML khác.
  \item Deep learning có thể rất chính xác nhưng \textbf{khó giải thích} hơn --- khó hiểu vì sao đưa ra kết luận.
\end{enumerate}
\end{ynghia}

\subsection{AI không nhất thiết qua ML}

\begin{luuy}[title={Triển khai AI bằng quy tắc}]
ML và deep learning đều là kỹ thuật giải bài toán AI.
Ta \textbf{có thể} triển khai AI trong thực tế mà không dùng ML/deep learning --- khi đó dùng thuật toán \textbf{dựa trên quy tắc} (rule-based).
\end{luuy}

\begin{dinhnghia}[title={Quan hệ tập hợp}]
Mọi phương pháp deep learning đều thuộc machine learning, nhưng \textbf{không phải} mọi machine learning đều là deep learning.
\end{dinhnghia}
